%! Author = sriadityavedantam
%! Date = 3/25/26

\lesson{10}{Friday 5 September 2025 9:10}{Day 10}

\begin{theorem}[Squeeze Theorem]{
    Suppose $(x_n),(y_n)$, and $(z_n)$ are sequences such that
    \[
        z_n\leq x_n\leq y_n \quad \forall n \geq N
    \]
    and $z_n, y_n \to L$.
    Then $\lim x_n = L$.
}
    Let $\epsilon > 0$ be given.
    There exists $N_1\in\n$ such that for all $n > N_1$, $\abs{y_n - L} < \epsilon$.
    There exists $N_2\in\n$ such that for all $n > N_2$, $\abs{z_n - L} < \epsilon$.
    Let $N \coloneqq \max\{N_1,N_2\}$.
    Then for all $n > N$, $y_n, z_n \in (L-\epsilon, L+\epsilon)$.
    Hence $[z_n, y_n] \subset (L-\epsilon, L+\epsilon)$.
    Since $x_n \in [z_n, y_n]$, we have $\abs{x_n - L} < \epsilon$.
\end{theorem}

\begin{theorem}{
    The sequence $(x_n)$ converges if and only if $\liminf x_n = \limsup x_n$.
}
    $(\impliedby)$.
    Suppose $\liminf x_n = \limsup x_n = L$.
    Then for all $n$, $\liminf x_n \leq x_n \leq \limsup x_n$.
    By the Squeeze Theorem, $x_n \to L$.
    $(\implies)$.
    Suppose $x_n \to L$.
    Let $\epsilon > 0$.
    Then there exists $N\in\n$ such that for all $n > N$, $\abs{x_n - L} < \epsilon$.
    Hence for all $n > N$, $L - \epsilon < x_n < L + \epsilon$.
    Thus for all $n > N$,
    \[
        \{x_m : m \geq n\} \subset (L - \epsilon, L + \epsilon).
    \]
    Let $y_n \coloneqq \sup\{x_m : m \geq n\}$ and $z_n \coloneqq \inf\{x_m : m \geq n\}$.
    Then for all $n > N$,
    \[
        L - \epsilon \leq z_n \leq y_n \leq L + \epsilon.
    \]
    Hence $\abs{y_n - L} < \epsilon$ and $\abs{z_n - L} < \epsilon$.
    Therefore $y_n \to L$ and $z_n \to L$.
    Thus $\limsup x_n = \liminf x_n = L$.
\end{theorem}

\begin{definition}[Diverges]
    A sequence $(x_n)$ diverges to $+\infty$ if for every $C\in\r$, there exists $N\in\n$ such that for all $n > N$, $x_n > C$.
    Similarly for $x_n \to -\infty$.
\end{definition}

\begin{lemma*}{
    Let $(x_n)$ be any sequence in $\r$.
    Suppose there is a sequence $(m_k)\subset\n$ such that $m_k \to +\infty$ and $(x_{m_k}) \to L$.
    Then there exists a subsequence of $(x_n)$ that converges to $L$.
}
\end{lemma*}

\begin{example}{
    For $n_k\in\n$ with $n_1 < n_2 < \ldots$, then $n_k \to +\infty$.
    Choose $n_k \geq k$ for all $k\in\n$.
}
\end{example}

\begin{proposition*}{
    $x_n \to \infty$ if and only if $x_{n_k} \to \infty$ for every subsequence $(x_{n_k})$.
}
\end{proposition*}